% XeLaTeX でコンパイルしてください
\documentclass[a4paper,11pt]{article}

% パッケージの読み込み
\usepackage{fontspec}
\setmainfont{Noto Sans CJK JP}[
  Script=CJK,
  Language=Japanese
]
\usepackage{amsmath, amssymb, amsthm}
\usepackage{tikz}
\usetikzlibrary{arrows.meta, positioning, calc, decorations.pathreplacing}
\usepackage{tcolorbox}
\tcbuselibrary{theorems}
\usepackage{hyperref}

% 定理環境の設定
\theoremstyle{definition}
\newtheorem{definition}{定義}[section]
\newtheorem{theorem}[definition]{定理}
\newtheorem{proposition}[definition]{命題}
\newtheorem{lemma}[definition]{補題}
\newtheorem{example}[definition]{例}
\newtheorem{remark}[definition]{注意}

% ボックス環境の設定
\newtcolorbox{keyidea}{
    colback=blue!5!white,
    colframe=blue!75!black,
    title=重要なアイディア
}

% タイトル情報
\title{構造塔の切断操作とその性質\\{\large Structure Tower の Slicing について}}
\author{%
    Lean 形式化: CODEX (OpenAI)\\
    \LaTeX 解説: Claude (Anthropic)
}
\date{\today}

\begin{document}

\maketitle

\begin{abstract}
本稿では、Bourbaki流の構造主義的数学における重要な概念である「構造塔」(Structure Tower)の切断操作について解説する。特に、与えられた構造塔を特定の添字より上で切り取る操作 \texttt{sliceAbove} の定義とその性質を、Lean 4 による形式化を通じて詳しく説明する。本解説は学部生を対象とし、直感的な理解を重視した説明を心がけた。
\end{abstract}

\tableofcontents

\section{導入}

\subsection{構造塔とは}

\begin{definition}[構造塔の直感的イメージ]
構造塔(Structure Tower)とは、数学的対象を「層」として積み重ねた構造である。各層には添字がついており、より豊かな構造を持つ層ほど「上」に位置する。
\end{definition}

\begin{example}
代数系の構造塔を考えてみよう:
\begin{itemize}
    \item 最下層:集合(set)
    \item 第2層:群構造を持つ集合(group)
    \item 第3層:環構造を持つ集合(ring)
    \item 第4層:体構造を持つ集合(field)
\end{itemize}
この例では、より高い層ほどより多くの公理を満たす「豊かな」構造を持つ。
\end{example}

\subsection{構造塔の可視化}

\begin{center}
\begin{tikzpicture}[scale=1.2]
    % 塔の描画
    \foreach \i/\name/\color in {0/集合/blue!20, 1/群/green!20, 2/環/orange!20, 3/体/red!20} {
        \fill[\color] (-2,\i) rectangle (2,\i+0.8);
        \draw[thick] (-2,\i) rectangle (2,\i+0.8);
        \node at (0,\i+0.4) {\name};
        \node[left] at (-2.2,\i+0.4) {$i_{\i}$};
    }
    
    % 要素の表示
    \foreach \i in {0,1,2,3} {
        \fill[black] (2.5,\i+0.4) circle (0.05);
        \fill[black] (3,\i+0.4) circle (0.05);
        \fill[black] (3.5,\i+0.4) circle (0.05);
    }
    
    % 矢印と説明
    \draw[-Stealth, thick] (4,2) -- (4,0.5);
    \node[right, align=left] at (4.2,1.25) {より豊かな\\構造へ};
    
    % タイトル
    \node[above] at (0,4) {\textbf{構造塔のイメージ}};
\end{tikzpicture}
\end{center}

\section{構造塔の形式的定義}

\subsection{構造塔の構成要素}

Lean 4 で定義される \texttt{StructureTowerWithMin} は以下の要素から構成される:

\begin{definition}[構造塔の構成要素]
構造塔 $T$ は以下のデータからなる:
\begin{enumerate}
    \item \textbf{台集合} (\texttt{carrier}): すべての要素を含む集合
    \item \textbf{添字集合} (\texttt{Index}): 層を区別するための添字の集合で、前順序が入っている
    \item \textbf{層関数} (\texttt{layer}): 各添字 $i$ に対して、その層に属する要素の集合 $\text{layer}(i)$ を対応させる
    \item \textbf{被覆性} (\texttt{covering}): すべての要素 $x$ は少なくとも一つの層に属する
    \item \textbf{単調性} (\texttt{monotone}): $i \leq j$ ならば $\text{layer}(i) \subseteq \text{layer}(j)$
    \item \textbf{最小層} (\texttt{minLayer}): 各要素 $x$ が属する最も低い層の添字
\end{enumerate}
\end{definition}

\begin{keyidea}
単調性により、下の層に属する要素は必ず上の層にも属する。これは「一度得られた構造は失われない」という直感に対応する。
\end{keyidea}

\section{切断操作: \texttt{sliceAbove}}

\subsection{切断操作の動機}

構造塔全体を扱う代わりに、「ある添字 $i$ より上の部分だけ」を取り出したい場合がある。これを実現するのが \texttt{sliceAbove} 操作である。

\begin{center}
\begin{tikzpicture}[scale=1.2]
    % 元の塔
    \begin{scope}[xshift=-4cm]
        \node[above] at (0,5) {\textbf{元の構造塔 $T$}};
        \foreach \i/\name in {0/層0, 1/層1, 2/層2, 3/層3, 4/層4} {
            \fill[blue!20] (-1.5,\i*0.9) rectangle (1.5,\i*0.9+0.7);
            \draw[thick] (-1.5,\i*0.9) rectangle (1.5,\i*0.9+0.7);
            \node at (0,\i*0.9+0.35) {\small \name};
        }
        
        % 切断線
        \draw[red, ultra thick, dashed] (-2,2*0.9-0.1) -- (2,2*0.9-0.1);
        \node[red, left] at (-2,2*0.9-0.1) {切断位置 $i$};
    \end{scope}
    
    % 矢印
    \draw[-Stealth, ultra thick] (-1.5,2.5) -- (1.5,2.5);
    \node[above] at (0,2.5) {\texttt{sliceAbove}$(T, i)$};
    
    % 切断後の塔
    \begin{scope}[xshift=4cm]
        \node[above] at (0,5) {\textbf{切断後 $T'$}};
        \foreach \i/\name in {2/層2, 3/層3, 4/層4} {
            \fill[green!20] (-1.5,\i*0.9) rectangle (1.5,\i*0.9+0.7);
            \draw[thick] (-1.5,\i*0.9) rectangle (1.5,\i*0.9+0.7);
            \node at (0,\i*0.9+0.35) {\small \name};
        }
        
        % 下の部分は透明に
        \foreach \i in {0,1} {
            \draw[dashed, gray] (-1.5,\i*0.9) rectangle (1.5,\i*0.9+0.7);
        }
        \node[gray] at (0,0.5) {\small (削除)};
    \end{scope}
\end{tikzpicture}
\end{center}

\subsection{形式的定義}

\begin{definition}[\texttt{sliceAbove} の定義]
構造塔 $T$ と添字 $i \in T.\text{Index}$ に対して、$\texttt{sliceAbove}(T, i)$ は以下のように定義される新しい構造塔 $T'$ である:

\begin{enumerate}
    \item \textbf{台集合}: 
    \[
    T'.\text{carrier} := \{x \in T.\text{carrier} \mid i \leq T.\text{minLayer}(x)\}
    \]
    つまり、最小層が $i$ 以上である要素のみを残す。
    
    \item \textbf{添字集合}:
    \[
    T'.\text{Index} := \{j \in T.\text{Index} \mid i \leq j\}
    \]
    添字も $i$ 以上のものだけを残す。
    
    \item \textbf{層関数}:
    \[
    T'.\text{layer}(j) := \{x \in T'.\text{carrier} \mid x \in T.\text{layer}(j)\}
    \]
    
    \item その他の性質(被覆性、単調性、最小層の性質)は元の塔 $T$ から継承される。
\end{enumerate}
\end{definition}

\subsection{重要な観察}

\begin{proposition}[切断の性質]
$T' = \texttt{sliceAbove}(T, i)$ とすると:
\begin{enumerate}
    \item $T'$ の任意の要素 $x$ について、$i \leq T.\text{minLayer}(x)$ が成り立つ
    \item $T'$ は $T$ の「部分構造塔」とみなせる
    \item 元の塔での最小層が $i$ 未満だった要素は完全に除外される
\end{enumerate}
\end{proposition}

\begin{proof}[証明のスケッチ]
(1) は台集合の定義から直接従う。

(2) $T'$ の各性質が $T$ の対応する性質から自然に導かれることを確認すればよい。特に:
\begin{itemize}
    \item 被覆性: 元の塔で要素 $x$ が層 $k$ に属していたとする($x \in T.\text{layer}(k)$)。$x \in T'.\text{carrier}$ なので $i \leq T.\text{minLayer}(x)$。最小層の最小性により $T.\text{minLayer}(x) \leq k$ なので $i \leq k$。よって $k \in T'.\text{Index}$ であり、$x \in T'.\text{layer}(k)$ が成り立つ。
    
    \item 単調性: $T$ の単調性がそのまま $T'$ に持ち込まれる。
\end{itemize}

(3) も台集合の定義から明らか。
\end{proof}

\section{具体例による理解}

\subsection{例:代数系の階層}

再び代数系の例を考えよう。

\begin{example}[環以上の構造だけを残す]
次のような構造塔を考える:
\begin{align*}
i_0 &: \text{集合} \\
i_1 &: \text{マグマ}(二項演算を持つ集合) \\
i_2 &: \text{半群}(結合律を満たすマグマ) \\
i_3 &: \text{モノイド}(単位元を持つ半群) \\
i_4 &: \text{群}(逆元を持つモノイド) \\
i_5 &: \text{可換群}(交換律を満たす群) \\
i_6 &: \text{環}(可換群に分配的な乗法を追加) \\
i_7 &: \text{可換環} \\
i_8 &: \text{整域} \\
i_9 &: \text{体}
\end{align*}

$\texttt{sliceAbove}(T, i_6)$ を適用すると、環以上の構造(環、可換環、整域、体)だけが残る。
\end{example}

\begin{center}
\begin{tikzpicture}[scale=0.9]
    % 元の塔(左側)
    \begin{scope}[xshift=-5cm]
        \node[above] at (0,6.5) {\small 元の塔};
        \foreach \i/\name/\col in {
            0/集合/blue!10,
            1/マグマ/blue!15,
            2/半群/blue!20,
            3/モノイド/blue!25,
            4/群/blue!30,
            5/可換群/blue!35,
            6/環/green!30,
            7/可換環/green!40,
            8/整域/green!50,
            9/体/green!60
        } {
            \fill[\col] (-1.2,\i*0.6) rectangle (1.2,\i*0.6+0.5);
            \draw[thick] (-1.2,\i*0.6) rectangle (1.2,\i*0.6+0.5);
            \node[font=\tiny] at (0,\i*0.6+0.25) {\name};
        }
        \draw[red, ultra thick, dashed] (-1.5,6*0.6-0.05) -- (1.5,6*0.6-0.05);
        \node[red, right, font=\tiny] at (1.5,6*0.6-0.05) {$i_6$};
    \end{scope}
    
    % 矢印
    \draw[-Stealth, thick] (-3,3) -- (-1,3);
    
    % 切断後の塔(右側)
    \begin{scope}[xshift=1cm]
        \node[above] at (0,6.5) {\small \texttt{sliceAbove}$(T,i_6)$};
        \foreach \i/\name/\col in {
            6/環/green!30,
            7/可換環/green!40,
            8/整域/green!50,
            9/体/green!60
        } {
            \fill[\col] (-1.2,\i*0.6) rectangle (1.2,\i*0.6+0.5);
            \draw[thick] (-1.2,\i*0.6) rectangle (1.2,\i*0.6+0.5);
            \node[font=\tiny] at (0,\i*0.6+0.25) {\name};
        }
        % 削除された部分
        \foreach \i in {0,...,5} {
            \draw[dashed, gray, opacity=0.3] (-1.2,\i*0.6) rectangle (1.2,\i*0.6+0.5);
        }
    \end{scope}
\end{tikzpicture}
\end{center}

\subsection{幾何学的解釈}

\begin{keyidea}
\texttt{sliceAbove} 操作は文字通り「塔を水平に切って上半分を取る」操作である。重要なのは:
\begin{itemize}
    \item 要素レベル: 最小層が $i$ 未満の要素は完全に除外される
    \item 層レベル: 添字が $i$ 未満の層は存在しなくなる
    \item 構造の保存: 残った部分は依然として構造塔の公理を満たす
\end{itemize}
\end{keyidea}

\section{Lean による形式化の詳細}

\subsection{型理論的な観点}

Lean コード中の重要な部分を見ていこう:

\begin{verbatim}
def sliceAbove (T : StructureTowerWithMin.{u, v}) 
    (i : T.Index) : StructureTowerWithMin.{u, v} where
  carrier := {x : T.carrier // i ≤ T.minLayer x}
  ...
\end{verbatim}

ここで \texttt{\{x : T.carrier // i ≤ T.minLayer x\}} は\textbf{部分型}(subtype)であり、条件 $i \leq T.\text{minLayer}(x)$ を満たす $T.\text{carrier}$ の要素全体を表す。

\subsection{証明の構造}

Lean コードでは、新しい構造塔 $T'$ が構造塔の公理を満たすことを証明する必要がある:

\begin{enumerate}
    \item \textbf{被覆性の証明}:
    \begin{verbatim}
covering := by
  intro x
  obtain ⟨k, hk⟩ := T.covering x.val
  have hik : i ≤ k := by calc ...
  use ⟨k, hik⟩
  exact hk
    \end{verbatim}
    
    任意の要素 $x \in T'.\text{carrier}$ に対して、元の塔での被覆性を使い、適切な層の存在を示す。
    
    \item \textbf{単調性の証明}:
    \begin{verbatim}
monotone := by
  intro j₁ j₂ hjj
  intro x hx
  exact T.monotone hjj hx
    \end{verbatim}
    
    元の塔の単調性をそのまま利用できる。
    
    \item \textbf{最小層の性質の証明}:
    
    最小層の定義、最小層が実際に要素を含むこと、最小性の3つの性質を証明する。
\end{enumerate}

\section{応用と発展}

\subsection{双対操作: \texttt{sliceBelow}}

\texttt{sliceAbove} の双対として、「添字 $i$ より下の部分だけを取る」操作 \texttt{sliceBelow} も定義できる。

\begin{center}
\begin{tikzpicture}[scale=1]
    \foreach \i/\name in {0/層0, 1/層1, 2/層2, 3/層3, 4/層4} {
        \pgfmathsetmacro{\opacity}{ifthenelse(\i>2,0.3,1)}
        \fill[blue!20, opacity=\opacity] (-1.5,\i*0.8) rectangle (1.5,\i*0.8+0.6);
        \draw[thick, opacity=\opacity] (-1.5,\i*0.8) rectangle (1.5,\i*0.8+0.6);
        \node[opacity=\opacity] at (0,\i*0.8+0.3) {\small \name};
    }
    \draw[red, ultra thick, dashed] (-2,3*0.8) -- (2,3*0.8);
    \node[red, right] at (2,3*0.8) {\texttt{sliceBelow}};
    \node[below] at (0,-0.5) {添字 $i$ 以下の部分を残す};
\end{tikzpicture}
\end{center}

\subsection{構造塔の準同型}

二つの構造塔 $T_1, T_2$ の間の準同型(structure-preserving map)を考えることができ、切断操作はこの準同型と可換性を持つ。

\subsection{圏論的視点}

構造塔全体は圏を成し、切断操作は特定の関手とみなせる。これは高度なトピックだが、構造塔の理論を体系的に理解する上で重要である。

\section{まとめ}

本稿では、構造塔の切断操作 \texttt{sliceAbove} について解説した。要点をまとめる:

\begin{enumerate}
    \item 構造塔は数学的構造を「層」として積み重ねたもの
    \item \texttt{sliceAbove} は塔を水平に切断し、上半分だけを取り出す操作
    \item 切断後も構造塔の公理が保たれることが証明できる
    \item この操作は型理論的に厳密に定式化でき、Lean で形式化可能
\end{enumerate}

構造塔の理論は Bourbaki の数学の基礎付けに源を発し、現代の形式化数学においても重要な役割を果たしている。\texttt{sliceAbove} のような基本操作を理解することで、より複雑な構造の操作や理論構築への道が開かれる。

\begin{thebibliography}{9}
\bibitem{bourbaki}
N. Bourbaki, \textit{Theory of Sets}, Springer, 2004.

\bibitem{lean}
The Lean Community, \textit{Theorem Proving in Lean 4}, 
\url{https://leanprover.github.io/theorem_proving_in_lean4/}

\bibitem{mathlib}
The mathlib Community, \textit{The Lean Mathematical Library},
\url{https://github.com/leanprover-community/mathlib4}
\end{thebibliography}

\appendix
\section{Lean コード全文}

参考のため、形式化の全体を掲載する:

\begin{verbatim}
import Mathlib.Data.Set.Basic
import Mathlib.Order.Basic
import MyProjects.ST.CAT2_complete

universe u v

namespace StructureTowerWithMin

variable {T : StructureTowerWithMin.{u, v}}

/-- 構造塔を添字 i より上で切断する操作
幾何学的には「塔を水平に切って上半分を取る」イメージ -/
def sliceAbove (T : StructureTowerWithMin.{u, v}) 
    (i : T.Index) : StructureTowerWithMin.{u, v} where
  carrier := {x : T.carrier // i ≤ T.minLayer x}
  Index := {j : T.Index | i ≤ j}
  indexPreorder := inferInstance
  layer j := {x : {x : T.carrier // i ≤ T.minLayer x} | 
              x.val ∈ T.layer j.val}
  covering := by
    intro x
    obtain ⟨k, hk⟩ := T.covering x.val
    have hik : i ≤ k := by
      calc
        i ≤ T.minLayer x.val := x.property
        _ ≤ k := T.minLayer_minimal x.val k hk
    use ⟨k, hik⟩
    exact hk
  monotone := by
    intro j₁ j₂ hjj
    intro x hx
    exact T.monotone hjj hx
  minLayer := fun x => ⟨T.minLayer x.val, x.property⟩
  minLayer_mem := by
    intro x
    exact T.minLayer_mem x.val
  minLayer_minimal := by
    intro x j hx
    exact T.minLayer_minimal x.val j.val hx

end StructureTowerWithMin
\end{verbatim}

\end{document}
